\documentclass[10pt,a4paper]{amsart}
\usepackage{amsmath,amssymb,amsthm,graphicx}
% Retained minimal preamble for the Stacks Project chapter varieties.tex.
% Verbatim excerpt of the chapter preamble: Xy-pic setup (native preamble lines 14-18), the native theorem
% declarations with the mathematical macros (native preamble lines 154-217), and the project own no-op
% definitions of the reference/slogan/history display environments (native wrapper.tex lines 1-3) that the
% counted statement uses for its slogan. Copied mechanically; nothing added, removed or redefined.
% For commutative diagrams we use Xy-pic
\usepackage[all]{xy}

% We use 2cell for 2-commutative diagrams.
\xyoption{2cell}
% Theorem environments.
%
\theoremstyle{plain}
\newtheorem{theorem}[subsection]{Theorem}
\newtheorem{proposition}[subsection]{Proposition}
\newtheorem{lemma}[subsection]{Lemma}

\theoremstyle{definition}
\newtheorem{definition}[subsection]{Definition}
\newtheorem{example}[subsection]{Example}
\newtheorem{exercise}[subsection]{Exercise}
\newtheorem{situation}[subsection]{Situation}

\theoremstyle{remark}
\newtheorem{remark}[subsection]{Remark}
\newtheorem{remarks}[subsection]{Remarks}

\numberwithin{equation}{subsection}

% Macros
%
\def\lim{\mathop{\mathrm{lim}}\nolimits}
\def\colim{\mathop{\mathrm{colim}}\nolimits}
\def\Spec{\mathop{\mathrm{Spec}}}
\def\Hom{\mathop{\mathrm{Hom}}\nolimits}
\def\Ext{\mathop{\mathrm{Ext}}\nolimits}
\def\SheafHom{\mathop{\mathcal{H}\!\mathit{om}}\nolimits}
\def\SheafExt{\mathop{\mathcal{E}\!\mathit{xt}}\nolimits}
\def\Sch{\mathit{Sch}}
\def\Mor{\mathop{\mathrm{Mor}}\nolimits}
\def\Ob{\mathop{\mathrm{Ob}}\nolimits}
\def\Sh{\mathop{\mathit{Sh}}\nolimits}
\def\NL{\mathop{N\!L}\nolimits}
\def\CH{\mathop{\mathrm{CH}}\nolimits}
\def\proetale{{pro\text{-}\acute{e}tale}}
\def\etale{{\acute{e}tale}}
\def\QCoh{\mathit{QCoh}}
\def\Ker{\mathop{\mathrm{Ker}}}
\def\Im{\mathop{\mathrm{Im}}}
\def\Coker{\mathop{\mathrm{Coker}}}
\def\Coim{\mathop{\mathrm{Coim}}}

% Boxtimes
%
\DeclareMathSymbol{\boxtimes}{\mathbin}{AMSa}{"02}

%
% Macros for moduli stacks/spaces
%
\def\QCohstack{\mathcal{QC}\!\mathit{oh}}
\def\Cohstack{\mathcal{C}\!\mathit{oh}}
\def\Spacesstack{\mathcal{S}\!\mathit{paces}}
\def\Quotfunctor{\mathrm{Quot}}
\def\Hilbfunctor{\mathrm{Hilb}}
\def\Curvesstack{\mathcal{C}\!\mathit{urves}}
\def\Polarizedstack{\mathcal{P}\!\mathit{olarized}}
\def\Complexesstack{\mathcal{C}\!\mathit{omplexes}}
% \Pic is the operator that assigns to X its picard group, usage \Pic(X)
% \Picardstack_{X/B} denotes the Picard stack of X over B
% \Picardfunctor_{X/B} denotes the Picard functor of X over B
\def\Pic{\mathop{\mathrm{Pic}}\nolimits}
\def\Picardstack{\mathcal{P}\!\mathit{ic}}
\def\Picardfunctor{\mathrm{Pic}}
\def\Deformationcategory{\mathcal{D}\!\mathit{ef}}
\newenvironment{reference}{}{}
\newenvironment{slogan}{}{}
\newenvironment{history}{}{}
\usepackage{cleveref}
\crefname{originallemma}{original lemma}{original lemmas}
\Crefname{originallemma}{Original lemma}{Original lemmas}
\makeatletter
\newlabel{algebra-lemma-NAK}{{00DV}{source}{}{sourcepointer-1}{}}
\newlabel{algebra-lemma-NAK@cref}{{[originallemma][0][]00DV}{[1][0][]source}{}{}{}}
\makeatother
\crefname{originallemma}{original lemma}{original lemmas}
\Crefname{originallemma}{Original lemma}{Original lemmas}
\makeatletter
\newlabel{algebra-lemma-unramified}{{00UV}{source}{}{sourcepointer-2}{}}
\newlabel{algebra-lemma-unramified@cref}{{[originallemma][0][]00UV}{[1][0][]source}{}{}{}}
\makeatother
\crefname{originallemma}{original lemma}{original lemmas}
\Crefname{originallemma}{Original lemma}{Original lemmas}
\makeatletter
\newlabel{lemma-tangent-space-cotangent-space}{{0B2D}{source}{}{sourcepointer-3}{}}
\newlabel{lemma-tangent-space-cotangent-space@cref}{{[originallemma][0][]0B2D}{[1][0][]source}{}{}{}}
\makeatother
\crefname{originallemma}{original lemma}{original lemmas}
\Crefname{originallemma}{Original lemma}{Original lemmas}
\makeatletter
\newlabel{morphisms-lemma-base-change-differentials}{{01V0}{source}{}{sourcepointer-4}{}}
\newlabel{morphisms-lemma-base-change-differentials@cref}{{[originallemma][0][]01V0}{[1][0][]source}{}{}{}}
\makeatother
\crefname{originallemma}{original lemma}{original lemmas}
\Crefname{originallemma}{Original lemma}{Original lemmas}
\makeatletter
\newlabel{morphisms-lemma-finite-type-differentials}{{01V2}{source}{}{sourcepointer-5}{}}
\newlabel{morphisms-lemma-finite-type-differentials@cref}{{[originallemma][0][]01V2}{[1][0][]source}{}{}{}}
\makeatother
\title{Factorization through affine space of relative dimension r with unramified projection at a point}
\author{The Stacks Project Authors}
\date{stacks-project:varieties.tex:local-snapshot-20261011; GNU FDL 1.2 or later}
\begin{document}
\maketitle
\noindent\textit{Editorial note.} All delivered mathematical text is the author's own, copied from the pinned Stacks Project chapter file varieties.tex; no statement or proof step was rewritten, reordered or removed.\par\medskip
\noindent\textit{Editorial note.} The proof is complete as written and ends with the Nakayama step; the criterion for unramifiedness is cited but the construction of the factorization is entirely inside the counted span.\par\medskip
\noindent\textit{Editorial note.} The counted spans carry 5 navigation stubs and no bibliography citation.\par\medskip
\noindent\textit{Editorial note.} The preamble is the chapter's retained minimal preamble (sources/stacks-varieties/retained-preamble.tex lines 1-77), a verbatim excerpt of the native preamble.\par\medskip
\noindent\textit{Editorial note.} The label lemma-unramfied-over-affine is defined only in varieties.tex, so no cross-chapter restatement had to be weighed for this card.\par\medskip
\noindent\textit{Allowed standard prerequisites (editorial).}\begin{itemize}
\item The module of differentials of a finite type morphism is finite and restricts to the fibre. Morphisms, Lemmas morphisms-lemma-finite-type-differentials and morphisms-lemma-base-change-differentials, invoked at proof lines 3488-3493. Supplies finiteness of the dimension r and the identification of the fibre of Omega\_\{X/S\} with Omega\_\{X\_s/s\}.
\item Tangent space as the dual of the fibre of the cotangent module. Same chapter, Lemma lemma-tangent-space-cotangent-space, invoked at proof line 3494. Supplies the identification of r with dim T\_\{X/S,x\}.
\item Criterion for unramifiedness by vanishing of the module of differentials. Algebra, Lemma algebra-lemma-unramified, invoked at proof lines 3508-3509. Supplies that Omega\_\{B/P,q\} = 0 is what has to be shown.
\item Nakayama's lemma for finite modules. Algebra, Lemma algebra-lemma-NAK, invoked at proof lines 3515-3516. Supplies the passage from the vanishing of the fibre of Omega\_\{B/P\} at q to the vanishing of Omega\_\{B/P,q\} itself.
\end{itemize}\medskip
\ifdefined\hypertarget\hypertarget{sourcepointer-1}{}\fi
\noindent\textit{Source reference.} Original lemma 00DV: sources/stacks-chow/tags.txt, source lines 512--512. Pointer outside the counted span, invoked once for context; the load-bearing argument of the counted statement is the author's proof reproduced in full in the proof span.\par\medskip
\ifdefined\hypertarget\hypertarget{sourcepointer-2}{}\fi
\noindent\textit{Source reference.} Original lemma 00UV: sources/stacks-chow/tags.txt, source lines 1125--1125. Pointer outside the counted span, invoked once for context; the load-bearing argument of the counted statement is the author's proof reproduced in full in the proof span.\par\medskip
\ifdefined\hypertarget\hypertarget{sourcepointer-3}{}\fi
\noindent\textit{Source reference.} Original lemma 0B2D: sources/stacks-varieties/pinned-original-source.tex, source lines 3014--3014. Pointer outside the counted span, invoked once for context; the load-bearing argument of the counted statement is the author's proof reproduced in full in the proof span.\par\medskip
\ifdefined\hypertarget\hypertarget{sourcepointer-4}{}\fi
\noindent\textit{Source reference.} Original lemma 01V0: sources/stacks-chow/tags.txt, source lines 2429--2429. Pointer outside the counted span, invoked once for context; the load-bearing argument of the counted statement is the author's proof reproduced in full in the proof span.\par\medskip
\ifdefined\hypertarget\hypertarget{sourcepointer-5}{}\fi
\noindent\textit{Source reference.} Original lemma 01V2: sources/stacks-chow/tags.txt, source lines 2431--2431. Pointer outside the counted span, invoked once for context; the load-bearing argument of the counted statement is the author's proof reproduced in full in the proof span.\par\medskip
\begin{lemma}
\label{lemma-unramfied-over-affine}
Let $f : X \to S$ be a finite type morphism of affine schemes.
Let $x \in X$ with image $s \in S$. Let
$$
r =
\dim_{\kappa(x)} \Omega_{X/S, x} \otimes_{\mathcal{O}_{X, x}} \kappa(x) =
\dim_{\kappa(x)} \Omega_{X_s/s, x} \otimes_{\mathcal{O}_{X_s, x}} \kappa(x) =
\dim_{\kappa(x)} T_{X/S, x}
$$
Then there exists a factorization
$$
X \xrightarrow{\pi} \mathbf{A}^r_S \to S
$$
of $f$ such that $\pi$ is unramified at $x$.
\end{lemma}

\begin{proof}
By Morphisms, Lemma \ref{morphisms-lemma-finite-type-differentials}
the first dimension is finite.
The first equality follows as the restriction of
$\Omega_{X/S}$ to the fibre is the module of differentials
from Morphisms, Lemma
\ref{morphisms-lemma-base-change-differentials}.
The last equality follows from Lemma \ref{lemma-tangent-space-cotangent-space}.
Thus we see that the statement makes sense.

\medskip\noindent
To prove the lemma write $S = \Spec(A)$ and $X = \Spec(B)$ and let
$A \to B$ be the ring map corresponding to $f$. Let $\mathfrak q \subset B$
be the prime ideal corresponding to $x$. Choose a surjection of $A$-algebras
$A[x_1, \ldots, x_t] \to B$. Since $\Omega_{B/A}$ is generated by
$\text{d}x_1, \ldots, \text{d}x_t$ we see that their images in
$\Omega_{X/S, x} \otimes_{\mathcal{O}_{X, x}} \kappa(x)$ generate
this as a $\kappa(x)$-vector space. After renumbering we may assume
that $\text{d}x_1, \ldots, \text{d}x_r$ map to a basis of
$\Omega_{X/S, x} \otimes_{\mathcal{O}_{X, x}} \kappa(x)$.
We claim that $P = A[x_1, \ldots, x_r] \to B$ is unramified at $\mathfrak q$.
To see this it suffices to show that $\Omega_{B/P, \mathfrak q} = 0$
(Algebra, Lemma \ref{algebra-lemma-unramified}).
Note that $\Omega_{B/P}$ is the quotient of $\Omega_{B/A}$ by the
submodule generated by $\text{d}x_1, \ldots, \text{d}x_r$.
Hence
$\Omega_{B/P, \mathfrak q} \otimes_{B_\mathfrak q} \kappa(\mathfrak q) = 0$
by our choice of $x_1, \ldots, x_r$.
By Nakayama's lemma, more precisely Algebra, Lemma \ref{algebra-lemma-NAK}
part (2) which applies as $\Omega_{B/P}$ is finite (see reference above),
we conclude that $\Omega_{B/P, \mathfrak q} = 0$.
\end{proof}

\end{document}
